% =============================================================================
\chapter{Architecture logicielle}
% =============================================================================

\section{Architecture hexagonale (Ports \& Adapters)}

Le backend suit une \textbf{architecture hexagonale modulaire}. Chaque module métier
(\texttt{auth}, \texttt{artwork}, \texttt{shop}, etc.) est découpé en trois couches :

\includediagram{hexagonal}{Architecture hexagonale — découpage en couches}

\section{Diagramme de packages}

La figure~\ref{fig:packages} illustre l'organisation modulaire du code source.

\includediagram[width=0.95\textwidth]{packages}{Diagramme de packages — structure modulaire}

\section{Structure interne d'un module type}

Chaque module sous \texttt{com.yowpainter.modules.*} respecte la même arborescence :

\begin{lstlisting}[language=]
com.yowpainter.modules.<module>/
  domain/
    model/          # Entites metier pures
    port/out/       # Interfaces RepositoryPort
  application/
    service/        # Cas d'utilisation (orchestration)
  infrastructure/
    adapter/
      in/web/       # Controllers REST + DTOs
      out/          # JPA, clients HTTP, etc.
\end{lstlisting}

\section{Diagramme de composants}

\includediagram{composants}{Diagramme de composants — flux de dépendances}

\section{Modules fonctionnels}

\begin{table}[h]
  \centering
  \caption{Inventaire des 14 modules métier}
  \label{tab:modules}
  \small
  \begin{tabular}{@{}llp{7cm}@{}}
    \toprule
    \textbf{Module} & \textbf{Service principal} & \textbf{Responsabilité} \\
    \midrule
    auth          & AuthService               & Authentification, profils, proxy Kernel auth \\
    artist        & ArtistService             & Profil artiste, slug, analytics \\
    artwork       & ArtworkService            & Œuvres, médias, likes, commentaires \\
    shop          & ShopService               & Produits, commandes, commerce Kernel \\
    event         & EventService              & Événements, réservations, billets QR \\
    payment       & PaymentService            & CamPay, wallet, retraits, polling \\
    subscription  & SubscriptionService       & Plans FREE/PRO/ELITE \\
    chat          & ChatMessageService        & Messagerie REST + WebSocket \\
    notification  & NotificationService       & Notifications via Kernel \\
    news          & NewsService               & Actualités artiste \\
    admin         & AdminService              & Administration plateforme \\
    file          & FileService               & Upload délégué Kernel \\
    search        & — (agrégation)            & Recherche globale multi-services \\
    system        & —                         & Health, version, metrics \\
    \bottomrule
  \end{tabular}
\end{table}

\section{Couche Kernel partagée}

Le package \texttt{com.yowpainter.shared.kernel} centralise l'intégration avec le
\kernel. Il expose 7 ports hexagonaux :

\begin{table}[h]
  \centering
  \caption{Ports Kernel (shared/kernel/port)}
  \label{tab:kernel-ports}
  \begin{tabular}{@{}ll@{}}
    \toprule
    \textbf{Port} & \textbf{Responsabilité} \\
    \midrule
    KernelAuthPort            & Sign-up, login, refresh, logout, reset password \\
    KernelOrganizationPort      & Création/approbation org, plans commerciaux \\
    KernelAdministrationPort  & Rôles, assignation utilisateurs \\
    KernelActorPort             & Onboarding acteurs BUSINESS/ART \\
    KernelProductPort           & CRUD produits commerce \\
    KernelSalesPort             & Commandes, confirm/cancel \\
    KernelFilePort              & Upload fichiers \\
    KernelNotificationPort      & Envoi/récupération notifications \\
    \bottomrule
  \end{tabular}
\end{table}
